\documentclass[10pt,a4paper]{amsart}
\usepackage[margin=19mm]{geometry}
\usepackage{amsmath,amssymb,mathtools}
\usepackage[scr=rsfs,cal=euler]{mathalfa}
\usepackage{xfrac}
\usepackage[unicode,hidelinks,bookmarks=false]{hyperref}
\newcommand{\ie}{i.e.\ }
\newcommand{\cf}{cf.\ }
\newcommand{\resp}{respectively\ }
\newcommand{\Subsection}{\subsection}
\providecommand{\nobreakdash}{\nobreak\hbox{-}}
\setlength{\emergencystretch}{2em}
\multlinegap0pt
\usepackage{bm}
\usepackage{bbm}
\usepackage{dsfont}
\usepackage{xspace}
\usepackage{cancel}
\usepackage{mathtools}
\usepackage{amsthm}
\makeatletter
\@ifundefined{theorem}{\newtheorem{theorem}{Theorem}}{}
\makeatother
\newcommand{\rpp}{r^{\prime\prime}}
\newcommand{\rp}{r^\prime}
\title[Exact Variance and Fano Factor for Arbitrary Level Crossings]{Exact Variance and Fano Factor for Arbitrary Level Crossings in Stationary Gaussian Processes}
\author{Shivang Rawat, Flaviano Morone, David J. Heeger, Stefano Martiniani}
\date{Source: arXiv:2605.25278v1 (CC BY 4.0)}
\begin{document}
\maketitle

\noindent\emph{Source and licence.} Exact Variance and Fano Factor for Arbitrary Level Crossings in Stationary Gaussian Processes. Version arXiv:2605.25278v1 (CC BY 4.0), distributed under the Creative Commons Attribution 4.0 International licence, as recorded in the arXiv licence metadata for this exact version and shown on the abstract page https://arxiv.org/abs/2605.25278v1. Full rights notice: licenses/arxiv-2605.25278-CC-BY-4.0.txt.\par\smallskip

\noindent\textit{Editorial scope.} Counted unit: the Theorem whose source label is \texttt{thrm:positivity\_of\_correlation\_term} in the article (source0.tex lines 1375--1383, source label \texttt{thrm:positivity\_of\_correlation\_term}), delivered together with the article's own complete original proof of it (source0.tex lines 1385--1473, one \texttt{proof} environment of 89 lines). No mathematics is written, repaired or re-derived: statement and proof are the article's own text line for line. Required context retained uncounted: none. Every definition, display and label of those lines is retained unchanged. Omitted from this package and not redistributed: 1--1374, 1384--1384, 1474--2020. Nothing inside a retained excerpt is omitted or altered: every retained body is delivered exactly as the article's lines are written, so no omission receipt of any kind accompanies this package. Citations: the retained excerpts carry no citation command at all, so no bibliography entry of the article is transcribed and no reference is redistributed. Reference adaptation: none was needed and none was made; every reference inside the delivered unit resolves inside it, so no unresolved reference or citation is produced. Printed slips kept literal: no printed word, symbol, hypothesis or number of the counted statement or of its proof was altered. Layout and class: the article's own class is \texttt{IEEEtran} with packages including graphicx, amsmath, amssymb, amsthm, mathtools, bm, xcolor, booktabs, cancel, soul, cite, tikz, epstopdf, hyperref; the wrapper is the lane's pinned \texttt{amsart} editorial wrapper at 10pt on A4, which restates the theorem-like environments the retained text needs and adds the article's own macro definitions that the retained text needs (\texttt{\textbackslash rp}, \texttt{\textbackslash rpp}), with extra wrapper packages (bm, bbm, dsfont, xspace, cancel, mathtools). The delivered numbering is the unit's own. Assets: no retained span contains a figure environment, a graphics-inclusion command, a PDF-page inclusion, a table or a TikZ picture; no image file is copied into this package, the package declares no asset dependency and no image file is redistributed with it. Licence: version arXiv:2605.25278v1 of this article is distributed under the Creative Commons Attribution 4.0 International licence, as recorded in the arXiv licence metadata for this exact version and shown on the abstract page https://arxiv.org/abs/2605.25278v1. A full rights notice is delivered at licenses/arxiv-2605.25278-CC-BY-4.0.txt. Characters: every retained excerpt is pure ASCII, so the delivered engine, XeLaTeX with its default Latin Modern text font, reports no missing character.
\begin{theorem} \label{thrm:positivity_of_correlation_term}
Let $r(t)$ be the autocorrelation function of a one-dimensional real Gaussian stationary process. For brevity, we introduce the following notations: $r(t)\rightarrow r$, $r(0)\rightarrow r_0$, $\rp(t) \rightarrow p$, $\rp(0)\rightarrow p_0$, $-\rpp(t) \rightarrow q$, and $-\rpp(0)\rightarrow q_0$. Then, for all $t \geq 0$ the following expressions are strictly positive,
\begin{equation}
\begin{split}
 &-\frac{r+r_0}{2\left(p^2+(q-q_0)(r+r_0)\right)} > 0 \\
 &-\frac{r_0-r}{2\left(p^2+(q+q_0)(r-r_0)\right)} > 0
\end{split}
\end{equation}
\end{theorem}

\begin{proof}
Since $r(t)$ is the autocorrelation function of a real stationary process, it has the spectral representation
\begin{equation}
    r(t)=\int_{-\infty}^{\infty}\cos(\omega t)\,dF(\omega),
\end{equation}
where $F$ is a nonnegative measure. Differentiating with respect to time, we obtain
\begin{equation}
    r'(t)=-\int_{-\infty}^{\infty}\omega\sin(\omega t)\,dF(\omega)
    \quad\text{and}\quad
    r''(t)=-\int_{-\infty}^{\infty}\omega^2\cos(\omega t)\,dF(\omega).
\end{equation}
In the notation defined in the theorem and using the integral formulae above, we have,
\begin{equation}
\begin{split}
    r &= \int_{-\infty}^{\infty}\cos(\omega t)\,dF(\omega) \\
    p &= -\int_{-\infty}^{\infty}\omega\sin(\omega t)\,dF(\omega) \\
    q &= \int_{-\infty}^{\infty}\omega^2\cos(\omega t)\,dF(\omega) \\
    r_0 &= \int_{-\infty}^{\infty}dF(\omega) \\
    p_0 &= 0 \\
    q_0 &= \int_{-\infty}^{\infty}\omega^2 \,dF(\omega)
\end{split}
\end{equation}
Using these relations, we have the following integral expressions for the relevant quantities,
\begin{equation}
    r+r_0=\int_{-\infty}^{\infty}\bigl(1+\cos(\omega t)\bigr)\,dF(\omega)
    \quad\text{and}\quad
    r_0-r=\int_{-\infty}^{\infty}\bigl(1-\cos(\omega t)\bigr)\,dF(\omega),
\end{equation}
and,
\begin{equation}
    q_0-q=\int_{-\infty}^{\infty}\omega^2\Bigl(1-\cos(\omega t)\Bigr)dF(\omega)
    \quad\text{and}\quad
    q_0+q=\int_{-\infty}^{\infty}\omega^2\Bigl(1+\cos(\omega t)\Bigr)dF(\omega).    
\end{equation}
Now, writing
\begin{equation}
    p^2=\left(\int_{-\infty}^{\infty}\omega\sin(\omega t)\,dF(\omega)\right)^2,    
\end{equation}
and applying the Cauchy--Schwarz inequality,
\begin{equation}
\left( \int_{S} f(x) g(x) \, d\mu(x) \right)^2 \le \left( \int_{S} f(x)^2 \, d\mu(x) \right) \left( \int_{S} g(x)^2 \, d\mu(x) \right), 
\end{equation}
gives us,
\begin{equation}
    p^2\le \left(\int_{-\infty}^{\infty}\omega^2\frac{\sin^2(\omega t)}{1+\cos(\omega t)}\,dF(\omega)\right)
    \left(\int_{-\infty}^{\infty}(1+\cos(\omega t))\,dF(\omega)\right).
\end{equation}
A short calculation shows that
\begin{equation}
    \frac{\sin^2(\omega t)}{1+\cos(\omega t)}
    =\frac{4\sin^2\bigl(\frac{\omega t}{2}\bigr)\cos^2\bigl(\frac{\omega t}{2}\bigr)}
    {2\cos^2\bigl(\frac{\omega t}{2}\bigr)}
    =2\sin^2\Bigl(\frac{\omega t}{2}\Bigr) = 1 - \cos\left(\omega t\right).
\end{equation}
Therefore,
\begin{equation}
    p^2\le \left(\int_{-\infty}^{\infty}\omega^2\left(1 - \cos\left(\omega t\right)\right)dF(\omega)\right)\left(\int_{-\infty}^{\infty}(1+\cos(\omega t))\,dF(\omega)\right),
\end{equation}
or,
\begin{equation}
    p^2 \leq \left(q_0-q\right) \left(r_0+r\right).
\end{equation}
Since equality in the Cauchy--Schwarz inequality can occur only if the integrands are proportional, which is excluded for a nondegenerate spectral measure, i.e., a spectral measure not concentrated on just a finite set of points, we have,
\begin{equation}
    p^2-(q_0-q)(r+r_0)<0.
\end{equation}
Further, since $r+r_0>0\; \forall \;t \geq 0$, we have,
\begin{equation}
 -\frac{r+r_0}{2\left(p^2+(q-q_0)(r+r_0)\right)} > 0 
\end{equation}
Applying Cauchy--Schwarz, we can also write,
\begin{equation}
\begin{split}
    p^2 &\leq \left(\int_{-\infty}^{\infty}\omega^2\frac{\sin^2(\omega t)}{1-\cos(\omega t)}\,dF(\omega)\right)
    \left(\int_{-\infty}^{\infty}(1-\cos(\omega t))\,dF(\omega)\right) \\
    p^2 &\leq \left(\int_{-\infty}^{\infty}\omega^2\left(1+\cos(\omega t)\right)\,dF(\omega)\right)
    \left(\int_{-\infty}^{\infty}(1-\cos(\omega t))\,dF(\omega)\right) \\
    p^2 &\leq \left(q_0+q\right)\left(r_0-r\right) \\
\end{split}
\end{equation}
Since the spectral measure is nondegenerate, strict inequality follows,
\begin{equation}
    p^2 -\left(q_0+q\right)\left(r_0-r\right) < 0
\end{equation}
Further, since $r_0-r>0\; \forall \;t \geq 0$, we have,
\begin{equation}
 -\frac{r_0-r}{2\left(p^2+(q+q_0)(r-r_0)\right)} > 0 
\end{equation}
\end{proof}

\end{document}
